\section{Torsion free modules}
\label{section-torsion-flat}

\noindent
In this section we discuss torsion free modules and the relationship
with flatness (especially over dimension 1 rings).

\begin{definition}
\label{definition-torsion}
Let $R$ be a domain. Let $M$ be an $R$-module.
\begin{enumerate}
\item We say an element $x \in M$ is {\it torsion} if there exists
a nonzero $f \in R$ such that $fx = 0$.
\item We say $M$ is {\it torsion free} if the only torsion element of $M$
is $0$.
\item We say $M$ is a {\it torsion module} if every element of $M$
is torsion.
\end{enumerate}
\end{definition}

\noindent
Let $R$ be a domain and let $S = R \setminus \{0\}$ be the multiplicative
set of nonzero elements of $R$. Then an $R$-module $M$ is torsion free if
and only if $M \to S^{-1}M$ is injective. In other words, if and only if
the map $M \to M \otimes_R K$ is injective where $K = S^{-1}R$ is the fraction
field of $R$.

\begin{lemma}
\label{lemma-torsion}
Let $R$ be a domain. Let $M$ be an $R$-module. The set of torsion elements
$M_{tors}$ of $M$ is the kernel of the map $M \to M \otimes_R K$.
Thus $M_{tors}$ is an $R$-submodule of $M$.
The quotient module $M/M_{tors}$ is torsion free.
\end{lemma}

\begin{proof}
See discussion above.
\end{proof}

\begin{lemma}
\label{lemma-localize-torsion}
Let $R$ be a domain. Let $M$ be a torsion free $R$-module.
For any multiplicative set $S \subset R$ the module
$S^{-1}M$ is a torsion free $S^{-1}R$-module.
\end{lemma}

\begin{proof}
Omitted.
\end{proof}

\begin{lemma}
\label{lemma-flat-pullback-torsion}
Let $R \to R'$ be a flat homomorphism of domains. If $M$ is a torsion
free $R$-module, then $M \otimes_R R'$ is a torsion free $R'$-module.
\end{lemma}

\begin{proof}
If $M$ is torsion free, then $M \subset M \otimes_R K$ is injective
where $K$ is the fraction field of $R$. Since $R'$ is flat over $R$
we see that $M \otimes_R R' \to (M \otimes_R K) \otimes_R R'$ is injective.
Since $M \otimes_R K$ is isomorphic to a direct sum of copies of $K$,
it suffices to see that $K \otimes_R R'$ is torsion free. This is true
because it is a localization of $R'$.
\end{proof}

\begin{lemma}
\label{lemma-extension-torsion-free}
Let $R$ be a domain. Let $0 \to M \to M' \to M'' \to 0$
be a short exact sequence of $R$-modules. If $M$ and $M''$
are torsion free, then $M'$ is torsion free.
\end{lemma}

\begin{proof}
Omitted.
\end{proof}

\begin{lemma}
\label{lemma-check-torsion}
Let $R$ be a domain. Let $M$ be an $R$-module.
Then $M$ is torsion free if and only if $M_\mathfrak m$ is a
torsion free $R_\mathfrak m$-module for all maximal ideals
$\mathfrak m$ of $R$.
\end{lemma}

\begin{proof}
Omitted. Hint: Use Lemma \ref{lemma-localize-torsion} and
Algebra, Lemma \ref{algebra-lemma-characterize-zero-local}.
\end{proof}

\begin{lemma}
\label{lemma-finite-torsion-free-submodule-free}
Let $R$ be a domain. Let $M$ be a finite $R$-module.
Then $M$ is torsion free if and only if $M$ is a
submodule of a finite free module.
\end{lemma}

\begin{proof}
If $M$ is a submodule of $R^{\oplus n}$, then $M$ is torsion free.
For the converse, assume $M$ is torsion free. Let $K$ be the
fraction field of $R$. Then $M \otimes_R K$ is a finite dimensional
$K$-vector space. Choose a basis $e_1, \ldots, e_r$ for this vector
space. Let $x_1, \ldots, x_n$ be generators of $M$. Write
$x_i = \sum (a_{ij}/b_{ij}) e_j$ for some $a_{ij}, b_{ij} \in R$
with $b_{ij} \not = 0$. Set $b = \prod_{i, j} b_{ij}$.
Since $M$ is torsion free the map
$M \to M \otimes_R K$ is injective and the image is contained
in $R^{\oplus r} = R e_1/b \oplus \ldots \oplus Re_r/b$.
\end{proof}

\begin{lemma}
\label{lemma-torsion-free-finite-noetherian-domain}
Let $R$ be a Noetherian domain. Let $M$ be a nonzero finite $R$-module.
The following are equivalent
\begin{enumerate}
\item $M$ is torsion free,
\item $M$ is a submodule of a finite free module,
\item $(0)$ is the only associated prime of $M$,
\item $(0)$ is in the support of $M$ and $M$ has property $(S_1)$, and
\item $(0)$ is in the support of $M$ and $M$ has no embedded associated prime.
\end{enumerate}
\end{lemma}

\begin{proof}
We have seen the equivalence of (1) and (2) in
Lemma \ref{lemma-finite-torsion-free-submodule-free}.
We have seen the equivalence of (4) and (5) in
Algebra, Lemma \ref{algebra-lemma-criterion-no-embedded-primes}.
The equivalence between (3) and (5) is immediate from the definition.
A localization of a torsion free module is torsion free
(Lemma \ref{lemma-localize-torsion}), hence it is clear that a
$M$ has no associated primes different from $(0)$. Thus (1)
implies (5). Conversely, assume (5). If $M$ has torsion,
then there exists an embedding $R/I \subset M$ for some nonzero
ideal $I$ of $R$. Hence $M$ has an associated prime different
from $(0)$
(see Algebra, Lemmas \ref{algebra-lemma-ass} and \ref{algebra-lemma-ass-zero}).
This is an embedded associated prime which contradicts the assumption.
\end{proof}

\begin{lemma}
\label{lemma-flat-torsion-free}
Let $R$ be a domain. Any flat $R$-module is torsion free.
\end{lemma}

\begin{proof}
If $x \in R$ is nonzero, then $x : R \to R$ is injective, and hence if $M$
is flat over $R$, then $x : M \to M$ is injective. Thus if $M$ is flat over
$R$, then $M$ is torsion free.
\end{proof}

\begin{lemma}
\label{lemma-valuation-ring-torsion-free-flat}
Let $A$ be a valuation ring.
An $A$-module $M$ is flat over $A$ if and only if $M$ is torsion free.
\end{lemma}

\begin{proof}
The implication ``flat $\Rightarrow$ torsion free'' is
Lemma \ref{lemma-flat-torsion-free}.
For the converse, assume $M$ is torsion free.
By the equational criterion of flatness (see
Algebra, Lemma \ref{algebra-lemma-flat-eq})
we have to show that every relation in $M$ is trivial. To do this assume that
$\sum_{i = 1, \ldots, n} a_i x_i = 0$ with $x_i \in M$ and $a_i \in A$.
After renumbering we may assume that $v(a_1) \leq v(a_i)$ for all $i$.
Hence we can write $a_i = a'_i a_1$ for some $a'_i \in A$. Note that
$a'_1 = 1$. As $M$ is torsion free we see that
$x_1 = - \sum_{i \geq 2} a'_i x_i$. Thus, if we choose
$y_i = x_i$, $i = 2, \ldots, n$ then
$$
x_1 = \sum\nolimits_{j \geq 2} -a'_j y_j, \quad
x_i = y_i, (i \geq 2)\quad
0 = a_1 \cdot (-a'_j) + a_j \cdot 1 (j \geq 2)
$$
shows that the relation was trivial (to be explicit the elements
$a_{ij}$ are defined by setting $a_{11} = 0$, $a_{1j} = -a'_j$
for $j > 1$, and $a_{ij} = \delta_{ij}$ for $i, j \geq 2$).
\end{proof}

\begin{lemma}
\label{lemma-dedekind-torsion-free-flat}
Let $A$ be a Dedekind domain (for example a discrete valuation ring
or more generally a PID).
\begin{enumerate}
\item An $A$-module is flat if and only if it is torsion free.
\item A finite torsion free $A$-module is finite locally free.
\item A finite torsion free $A$-module is finite free if
$A$ is a PID.
\end{enumerate}
\end{lemma}

\begin{proof}
(For the parenthetical remark in the statement of the lemma, see
Algebra, Lemma \ref{algebra-lemma-PID-dedekind}.) Proof of (1).
By Lemma \ref{lemma-check-torsion}
and Algebra, Lemma \ref{algebra-lemma-flat-localization}
it suffices to check the statement over $A_\mathfrak m$
for $\mathfrak m \subset A$ maximal. Since $A_\mathfrak m$
is a discrete valuation ring
(Algebra, Lemma \ref{algebra-lemma-characterize-Dedekind})
we win by Lemma \ref{lemma-valuation-ring-torsion-free-flat}.

\medskip\noindent
Proof of (2). Follows from
Algebra, Lemma \ref{algebra-lemma-finite-projective}
and (1).

\medskip\noindent
Proof of (3). Let $A$ be a PID and let $M$ be a finite
torsion free module. By Lemma \ref{lemma-finite-torsion-free-submodule-free}
we see that $M \subset A^{\oplus n}$ for some $n$. We argue that
$M$ is free by induction on $n$. The case $n = 1$ expresses exactly the
fact that $A$ is a PID. If $n > 1$ let $M' \subset A^{\oplus n - 1}$
be the image of the projection onto the last $n - 1$ summands of
$A^{\oplus n}$. Then we obtain a short exact sequence
$0 \to I \to M \to M' \to 0$ where $I$ is the intersection of $M$
with the first summand $A$ of $A^{\oplus n}$. By induction we
see that $M$ is an extension of finite free $A$-modules, whence
finite free.
\end{proof}

\begin{lemma}
\label{lemma-hom-into-torsion-free}
Let $R$ be a domain. Let $M$, $N$ be $R$-modules.
If $N$ is torsion free, so is $\Hom_R(M, N)$.
\end{lemma}

\begin{proof}
Choose a surjection $\bigoplus_{i \in I} R \to M$.
Then $\Hom_R(M, N) \subset \prod_{i \in I} N$.
\end{proof}




